\subsection{Character of the Representation}

\[
    \chi_h(q) = \text{Tr}_{\mathrm{Vir}_h(c)} q^{L_0 - \frac{c}{24}},
\]
\(q\) is a formal variable, around which we can expand the character as a power series. The coefficient of \(q^n\) counts the number of states at level \(n\) in the representation of \(\mathrm{Vir}_h(c)\). The xponent has physical meaning: \(L_0\) is the dilatation operator, and the shift by \(-\frac{c}{24}\) is related to the Casimir energy of the vacuum state in a conformal field theory.

If we call \textit{verma} the representation generated by acting with the negative modes of the Virasoro algebra on the highest weight state \(\ket{h}\) and without null vectors, then we can write
\[
    \mathrm{Verma} = \frac{q^{h - \frac{c}{24}}}{\prod_{n=1}^\infty (1 - q^n)}.
\]
this is the euler function, which counts the number of ways to distribute the negative modes \(L_{-n}\) among the states in the Verma module. The factor \(q^{h - \frac{c}{24}}\) accounts for the conformal weight of the highest weight state and the Casimir energy shift. The denominator is a product over all positive integers \(n\), reflecting the fact that we can have any number of negative modes acting on the highest weight state to generate new states in the representation. If we take again the expression for the character of the irreducible representation, we can write
\[
    \begin{aligned}
        \chi_h(q) = \sum_{\ket{s} \in \mathrm{Vir}_h(c)} \bra{s} q^{L_0 - \frac{c}{24}} \ket{s} = q^{-\frac{c}{24}}\sum_{\ket{s} \in \mathrm{Vir}_h(c)} q^{L_0} \\
        = q^{-\frac{c}{24}} \sum_{n=0}^\infty \sum_{\ket{s} \in \mathrm{Vir}_h(c)} q^{h+n}\bra{s} \ket{s} = q^{h - \frac{c}{24}} \sum_{n=0}^\infty \left( \mathrm{dim} \mathcal{V}_{c,h} \right) q^n = q^{h - \frac{c}{24}} \sum_{n} a_n q^n,
    \end{aligned}
\]
where \(a_n\) counts the number of states at level \(n\) in the irreducible representation. Comparing this expression with the character of the Verma module, we can see that the presence of null vectors reduces the number of states at each level, leading to a different counting of states in the irreducible representation compared to the Verma module.

    [...]

\subsubsection{In Minimal Models}

In minimal models we have
\[
    \mathrm{Minimal} = \frac{q^{-\frac{c}{24}}}{\prod_{n=1}^\infty (1 - q^n)} \sum_{n \in \mathbb{Z}} q^{\frac{[2m(m+1)n + (m+1)r - ms]^2 - 1}{4m(m+1)}} - q^{\frac{s \to -s}{4m(m+1)}},
\]
where the sum over \(n\) accounts for the contributions of the null vectors in the representation, and the subtraction of the term with \(s \to -s\) ensures that we are counting only the states in the irreducible representation. The character of the minimal model can be expressed in terms of the characters of the Verma modules, reflecting the structure of the irreducible representations in terms of the Verma modules and their null vectors.

When we want to subtract the contribution of the null vectors, we have also to consider all the descendants of the null vectors, which are also null. This is reflected in the structure of the character, where the contributions of the null vectors and their descendants are subtracted from the character of the Verma module to obtain the character of the irreducible representation in the minimal model.

    [\textit{picture of nested cones representing verma modules and null vectors descendant}]

The main contribution to the character of the minimal model came from Kaz, but this last general expression was derived by Rocha-Caridi, one of his students.

Let us say that an important aspect of this topic can be understood only knowing the theory of theta functions, which we will not cover in these notes. Also ramanujan and hardi have contributed to the development of theta funxtion theory, and their work has been crucial in understanding the structure of characters in conformal field theories, especially in the context of minimal models. But these functions had some counting properties, which are the ones giving us the previous results, and which are related to the counting of states in the representations of the Virasoro algebra.

the coefficients in the expansion of q are the one telling us how many states we have at each level, without the null vectors. The presence of null vectors reduces the number of states at each level, leading to a different counting of states in the irreducible representation compared to the Verma module. For example at level 2 where we have \(L_{-2}\) and \(L_{-1}^2\), we have a null state in correspondence of the minimal models, and the coefficient of \(q^2\) in the character of the minimal model will be 1 instead of 2, reflecting the fact that one of the states at level 2 is null and does not contribute to the counting of states in the irreducible representation.

\subsubsection{Fusion rules}

Let us consider the most general fusion rule among two primary fields \(\phi_{r_1, s_1}\) and \(\phi_{r_2, s_2}\) in a minimal model (around a kaz table)
\[
    \phi_{r_1, s_1} \times \phi_{r_2, s_2} = \sum_{r=|r_1 - r_2| + 1,\; r = r_1 + r_2 + 1 (\mathrm{mod} 2)}^{\min(r_1 + r_2 - 1, 2m - r_1 - r_2 - 1)} \sum_{s=|s_1 - s_2| + 1\dots }^{\min(s_1 + s_2 - 1, 2(m+1) - s_1 - s_2 - 1)} \phi_{r,s},
\]
which seems to be a very complicated expression, but it is actually quite simple. The fusion rules in minimal models are determined by the structure of the representations of the Virasoro algebra and the presence of null vectors. The indices \(r\) and \(s\) in the primary fields \(\phi_{r,s}\) label the representations of the Virasoro algebra, and the fusion rules dictate how these representations combine when we take the product of two primary fields.

The sum is very similar to the composition of \(\mathrm{SU}(2)\) representations: we are summing from the absolute value of the difference of the indices plus one, up to the minimum of the sum of the indices minus one and a certain cutoff determined by the parameters of the minimal model. The conditions on \(r\) and \(s\) ensure that we are only considering valid representations that can arise from the fusion of the two primary fields. We are summing in modulo 2 since we are only considering integer values of \(r\) and \(s\) that satisfy the fusion rules. Let us write a couple of examples to see how this works in practice. For example, if we take the fusion of \(\phi_{1,2}\) with himself, we have
\[
    \phi_{1,2} \times \phi_{1,2} = \phi_{1,1} + \phi_{1,3},
\]
where we have used the fusion rules to determine which primary fields appear in the product. In this case, the fusion of \(\phi_{1,2}\) with itself yields two primary fields: \(\phi_{1,1}\) and \(\phi_{1,3}\).

In principle we can write the katz table from these compositions, and it would also be infinite, but at a certain point, due to the symmetry, at a certain point we reach the \textit{truncation from above}: we will find an identity operator along the family of descendants of the null vector, and this will lead to a truncation of the representations, which is one of the key features of minimal models. Thus we have a way to compute the Katz table from the fusion rules, and from here we can build up the whole structure of the minimal model, we can write complete OPEs and correlation functions.

Futhermore, the presence of a null vector implies a differential equation for the correlation functions, which can be used to compute the correlation functions explicitly. The fusion rules also play a crucial role in determining the structure of the operator product expansion (OPE) in minimal models, as they dictate which primary fields can appear in the OPE of two given primary fields. This allows us to compute correlation functions and understand the structure of the theory in a systematic way.




\section{modular invariance}

we are going to show a method which allows us to classify all the minimal models, which is one of the most fruitful approach to CFT. We will be able to distinguish relevant and irrelevant operators, and to compute the critical exponents of the theory.

in string theory, if we impose modular invariance fora string in a wordsheet it gives us its unitarity...

The idea is that if the theroy can be put on a plane and on a cilinder and we have a conformal transformation which maps the plane to the cilinder. Cardi also mapped theories from the plane to the half plane (where we have a boundary) but we can also write theories on annulus or disks, whichever its the more convenient choice. This will have consequences also on the quantum hole effect which we are not going to discuss.

\subsubsection{Geometry of the torus}

We can write the theory on the riemann sphere (complex plane) which is a surface of zero genus (we can deform any closed path on the sphere to a point continuously). On the thorus its the same, but if we consider a path which goes around the hole of the torus, we cannot deform it to a point, so the torus is not topologically equivalent to the sphere (we do not have diffeomorphisms which can map one to the other).

    [\textit{drawing of torus with path not shrinkable to a point}]

    [...]

First write the theory on a cylinder, which is topologically equivalent to the plane, and then we can map the cylinder to the torus by identifying the two ends of the cylinder, this is the idea. We have periodic boundary conditions in the spatial direction, but we can impose it also in time: let us choose a starting time at \(r=0\) and a final time at \(r=\tau\), and then we can identify these two times, so we have periodic boundary conditions also in time. This is the idea of writing the theory on a torus.

\begin{example}[Ising model as 2dim lattice in a computer]
    If we want to model the Ising model on a computer, we can write it on a torus, which is equivalent to imposing periodic boundary conditions in both spatial and temporal directions.

    Indeed we do not have an infinite memroy computer, so we cannot write the theory on an infinite lattice, but we can write it on a finite lattice with periodic boundary conditions, which is equivalent to writing the theory on a torus. This allows us to study the properties of the Ising model in a finite volume, and to extract information about the critical behavior of the model by analyzing the finite-size scaling of observables as we approach the critical point.

    The ability to write the theory on a torus and to impose periodic boundary conditions is crucial for numerical simulations of the Ising model and other lattice models, as it allows us to study the behavior of the system in a finite volume while still capturing the essential features of the critical behavior.
\end{example}

But the russians were matematicians, not computer scientist (fatteyev had difficulties to write an email), so they were not interested in numerical simulations, but they were interested in the mathematical structure of the theory, and in particular in the modular invariance of the theory on the torus, which is a powerful tool to classify the minimal models and to understand their properties.

In order to map a plane to the torus, we need to identify points in the complex plane in order to obtain doubly periodic functions, which are the functions defined on the torus (theta functions indeed are doubly periodic).

    [\textit{drawing of complex plane, parallelogram identified by two vectors and periodic conditions in the two dimensions}]

We can define a lattice as a linear combination of two vectors \(\omega_1\) and \(\omega_2\), which are not collinear, and we can identify points in the complex plane that differ by a lattice vector, so we have an equivalence relation on the complex plane given by \(z \sim z + m\omega_1 + n\omega_2\) for integers \(m\) and \(n\). This allows us to define a torus as the quotient of the complex plane by this lattice
\[
    \mathbb{T} = \mathbb{C} / \Lambda, \quad \text{where} \quad \Lambda = \{ m\omega_1 + n\omega_2 \mid m,n \in \mathbb{Z} \}.
\]
but we can do better: since we are in conformal invariance, the torus depend only on the ratio of the two vectors, so we can define a complex parameter
\[
    \tau = \frac{\omega_2}{\omega_1},
\]
which is called the modular parameter of the torus. We can always rotate one of the two vectors to be along the real axis, so we can choose \(\omega_1\) to be real and positive, and then \(\tau\) will be a complex number in the upper half-plane, since we need \(\mathrm{Im}(\tau) > 0\) to ensure that the torus is well-defined.

\(\tau\) is called \textbf{modular parameter} of the torus, and it encodes the shape of the torus. Different values of \(\tau\) correspond to different shapes of the torus, and the modular invariance of the theory on the torus means that the physical observables should be invariant under transformations of \(\tau\) that correspond to different ways of describing the same torus.

Topologically all the tori are the same, since we can map one to the other by a diffeomorphism, but we have different conformal structures on the torus, which are encoded in the modular parameter \(\tau\): if we consider only conformal transformations, we cannot map a torus with a given \(\tau\) to a torus with a different \(\tau\).

There is a but: if we consider a vector \((\omega_1,\, \omega_2)\), we can describe it from a basis on the lattice to another basis
\[
    \begin{pmatrix}
        \omega_1' \\
        \omega_2'
    \end{pmatrix} = \begin{pmatrix}
        a & b \\
        c & d
    \end{pmatrix} \begin{pmatrix}
        \omega_1 \\
        \omega_2
    \end{pmatrix},
\]
where \(a\), \(b\), \(c\), and \(d\) are integers, and the determinant of the matrix is 1 (after some renormalization), so we have \(ad - bc = 1\).\footnote{This transformation is in the special linear group \(\mathrm{SL}(2, \mathbb{Z})\), where only integer entries are allowed.} This transformation corresponds to a change of basis in the lattice, and it leaves the lattice invariant, since we are still describing the same set of points in the complex plane. However, this transformation changes the modular parameter \(\tau\) according to the formula
\[
    \tau' = \frac{a\tau + b}{c\tau + d}.
\]
So the value of the modular parameter \(\tau\) changes, but it remains the same if we consider the transformation matrix with a minus sign, since the determinant is 1, so we can consider transformations in the projective special linear group \(\mathrm{PSL}(2, \mathbb{Z})\), which is the quotient of \(\mathrm{SL}(2, \mathbb{Z})\) by its center. The modular invariance of the theory on the torus means that the physical observables should be invariant under these transformations of \(\tau\), which correspond to different ways of describing the same torus.

So we consider the complex plane of \(\tau\) and we have that a torus with a given \(\tau\) is equivalent to a torus with a different \(\tau'\) if they are related by a transformation in \(\mathrm{PSL}(2, \mathbb{Z})\). Thus we can select a region out of which the values of tau are equivalent to what we already found within the region

    [\textit{drawing of semicircles with center on the real axis and repetitions}]

All elements \(A \in \mathrm{PSL}(2, \mathbb{Z})\) can be generated by two transformations, called \(S\) and \(T\), by composing them in different ways:
\[
    A = S T T S S S T\dots
\]
and these transformations have the following matrix representation in \(\mathrm{PSL}(2, \mathbb{Z})\):
\[
    T = \begin{pmatrix}
        1 & 0 \\
        1 & 1
    \end{pmatrix}, \quad S = \begin{pmatrix}
        0 & -1 \\
        1 & 0
    \end{pmatrix},
\]
where
\[
    S^2 = \mathbb{I}, \quad (ST)^3 = \mathbb{I},
\]
which are the two constraints that the generators must satisfy. So if we want to impose the invariance under the full group \(\mathrm{PSL}(2, \mathbb{Z})\), we can just impose the invariance under the generators \(S\) and \(T\), which is a much simpler task. The transformation \(T\) corresponds to a Dehn twist of the torus, while the transformation \(S\) corresponds to an exchange of the two cycles of the torus:
\[
    T \,:\, \tau \mapsto \tau + 1, \quad S \,:\, \tau \mapsto -\frac{1}{\tau}.
\]
Indeed this group is called the \textbf{modular group}, and if we prove something to be invariant under the generators \(S\) and \(T\), we have that it is invariant under the full modular group, thus it has \textbf{modular invariance} and it will be invariant under all transformations that correspond to different ways of describing the same torus.

If we can implement the partition function on a torus in such a way that is invariant under the generators \(S\) and \(T\), we have that it is invariant under the full modular group, and this will allow us to classify the minimal models and to understand their properties. The partition function is a physical quantity that should be invariant under the modular transformations, and it encodes information about the spectrum of the theory and the structure of the representations of the Virasoro algebra.

\subsubsection{Partition Function and Modular Invariance}

Firstly we pass from the plane to the cylinder by the conformal transformation
\[
    z = e^{\frac{2\pi i}{\omega_1}w},
\]
where \(\omega_1\) is the first vector of the lattice, and \(w\) is the coordinate on the cylinder. This transformation maps the complex plane to a cylinder of circumference \(\omega_1\). We have a Hamiltonian of the form
\[
    H = \frac{2\pi}{\omega_1} \left( L_0 + \bar{L}_0 - \frac{c}{12}\right),
\]
while the momentum operator is given by
\[
    P = \frac{2\pi}{\omega_1} \left( L_0 - \bar{L}_0 \right).
\]

We can now compute the transfer matrix, which is given by
\[
    \hat{T} = e^{-H \Im(\omega_2) + i P \Re(\omega_2)},
\]
where we took the imaginary part of \(\omega_2\) since the displacements on the torus are given by the part of \(\omega_2\) that is orthogonal to \(\omega_1\), and we took the real part of \(\omega_2\) as a phase factor since it corresponds to a shift in the spatial direction. We can develop this expression using the definitions of \(H\) and \(P\) to obtain
\[
    \hat{T} = e^{-i \pi \frac{\omega_2 - \bar{\omega}_2}{\omega_1} \left( L_0 + \bar{L}_0 - \frac{c}{12}\right) + i\pi \frac{\omega_2 + \bar{\omega}_2}{\omega_1} \left( L_0 - \bar{L}_0 \right)},
\]
which with some algebraic manipulations can be rewritten as
\[
    \hat{T} = e^{2\pi i \tau (L_0 - \frac{c}{24})} e^{-2\pi i \bar{\tau} (\bar{L}_0 - \frac{c}{24})} = q^{L_0 - \frac{c}{24}} \bar{q}^{\bar{L}_0 - \frac{c}{24}},
\]